% ----------------------------------------------------------
% Introdução
% ----------------------------------------------------------
\chapter[Introdução]{Introdução}

No cenário contemporâneo da indústria de software, a capacidade de entregar funcionalidades de forma rápida, confiável e frequente consolidou-se como um diferencial competitivo crítico. O modelo tradicional de desenvolvimento, historicamente caracterizado por ciclos de lançamento longos e transferências manuais entre equipes de desenvolvimento e operações, mostra-se cada vez mais inadequado para atender à demanda atual por agilidade e inovação contínua. Esse cenário de pressão por velocidade frequentemente expõe gargalos operacionais, onde a tentativa de acelerar entregas sem a devida automação resulta em processos manuais propensos a erros humanos e instabilidade em produção, gerando um ambiente de tensão e risco elevado a cada novo lançamento.

Em resposta a esses desafios operacionais, a indústria tem adotado abordagens que visam unificar o desenvolvimento e as operações através de uma cultura de colaboração e, fundamentalmente, através da automação de processos. No coração técnico dessa transformação estão as práticas de Integração Contínua e Entrega Contínua, orquestradas por pipelines de implantação automatizados. Embora os benefícios teóricos dessas práticas, como a redução de riscos e o aumento da qualidade, sejam amplamente reconhecidos, muitas organizações ainda encontram dificuldades para quantificar objetivamente o impacto real dessa transição tecnológica. Apenas recentemente o mercado passou a adotar métricas mais científicas de desempenho para desafiar a antiga crença de que a velocidade de entrega deve ocorrer necessariamente em detrimento da estabilidade do sistema.

Este Trabalho de Conclusão de Curso tem como objetivo analisar quantitativamente o impacto da implementação de um pipeline de CI/CD na eficiência da entrega de software. Através de um estudo experimental comparativo, utilizando aplicações web modernas em diferentes pilhas tecnológicas e um ambiente de nuvem emulado localmente para garantir a reprodutibilidade, este trabalho irá mensurar a diferença de desempenho entre um processo de deploy manual e um processo totalmente automatizado. A análise será fundamentada em métricas de tempo de entrega e frequência de implantação, buscando validar empiricamente como a automação atua como um catalisador para a alta performance na engenharia de software.
